\section{Major arcs}\label{sec:main}

Take $C\ge3$ and set
\begin{equation}\label{eq:N}
 N=\left\lceil\frac C{\sctrl}\right\rceil.
\end{equation}
Choose $k_0$ sufficiently large that
\begin{equation}\label{eq:main-geometry}
 2N<2^{2k_0},
 \qquad 2N+1\le L,
 \qquad \frac1{\sE}\le N,
\end{equation}
and so that the Taylor estimates below hold.

For $|m|\le N$, put
\[
 T_m=
 \prod_{e\in E}
 \bigl((1-\theta)+\theta e^{2\pi im/e}\bigr)
 e^{-2\pi im/b}.
\]
For real $z$ define
$B_\theta(z)=(1-\theta)+\theta e^{2\pi iz}$.  If $|z|\le1/10$ and
$1/3\le\theta\le2/3$, then
\[
 \Re B_\theta(z)=1-\theta+\theta\cos(2\pi z)>0.
\]
Thus every factor lies in the open right half-plane and is nonzero, and the
principal logarithm is defined continuously along the segment from $0$ to
$z$.  On this disk,
\begin{equation}\label{eq:log-expansion}
 \log B_\theta(z)
 =2\pi i\theta z
  -2\pi^2\theta(1-\theta)z^2+R_\theta(z),
 \qquad |R_\theta(z)|\le C_3|z|^3.
\end{equation}
After a further increase of $k_0$, for every $e\in E$ and $|m|\le N$,
\begin{equation}\label{eq:Taylor-fields}
 |m/e|\le\frac1{10},
 \qquad
 C_3\sum_{e\in E}|m/e|^3\le\frac1{10}.
\end{equation}
These estimates follow from the denominator scales and
\eqref{eq:inverse-square-constant}: the smallest denominator tends to infinity
with $k_0$, while the inverse-square sum is bounded by a fixed multiple of
$\sctrl^2$, so the cubic remainder is uniformly $o(1)$ on this window.

Summing \eqref{eq:log-expansion}, the linear term cancels by
\eqref{eq:theta-window}, and
\[
 \log T_m=-2\pi^2m^2\sE^2+\delta_m,
 \qquad |\delta_m|\le\frac1{10}.
\]
Hence
\begin{equation}\label{eq:per-main}
 \Re T_m\ge0.8e^{-2\pi^2m^2\sE^2}.
\end{equation}
Conjugation sends $T_m$ to $T_{-m}$, so the major-arc sum is real.  The
integers $|m|\le1/(2\sE)$ lie in the major-arc window and each contributes at
least $0.8e^{-\pi^2/2}$.  A coarse integer count gives
\begin{equation}\label{eq:main-lower}
 \Re\sum_{|m|\le N}T_m
 \ge\frac{c_3}{\sE},
 \qquad
 c_3=0.8\frac{e^{-\pi^2/2}}2>0.
\end{equation}

For every real $t$,
\[
 |(1-\theta)+\theta e^{2\pi it}|^2
 =1-4\theta(1-\theta)\sin^2(\pi t).
\]
On the centered interval,
$\sin^2(\pi t)\ge4\normdist t^2$, and
$\theta(1-\theta)\ge2/9$.  Therefore
\begin{equation}\label{eq:fourier-energy}
 |F(h)|\le
 \exp\!\left(-\frac{16}{9}\QE(h)\right).
\end{equation}
